\chapter{Installazione e ambiente Conda}
\label{chap:conda}

\section{Architettura degli ambienti}

Gli entrypoint nella root, come \filepath{run.py} ed
\filepath{evaluate.py}, devono essere eseguiti nell'ambiente Conda principale
\code{graph_rec}, descritto da \filepath{environment.yml}. L'interprete usato
per VGGT è indicato da \envvar{VGGT_EMBEDDED_PYTHON} e può coincidere con quello
di \code{graph_rec}; se la variabile non è impostata viene usato l'interprete
dell'ambiente attivo.

La pipeline si avvia quindi dal normale ambiente Conda. Un runtime separato,
quando necessario, va indicato esplicitamente nella configurazione locale.

\manualnote[Due file diversi]{\filepath{environment.yml} descrive Python e
pacchetti dell'ambiente Conda; \filepath{.env} contiene invece percorsi e
opzioni della pipeline.}

\subsection{Creazione dell'ambiente}

Aprire un terminale nel quale il comando \code{conda} sia disponibile,
spostarsi nella root del repository e creare l'ambiente dichiarato dal progetto:

\begin{lstlisting}[language={}]
conda env create -f environment.yml
conda activate graph_rec
\end{lstlisting}

Per riallineare un ambiente già esistente al file versionato:

\begin{lstlisting}[language={}]
conda env update -n graph_rec -f environment.yml --prune
\end{lstlisting}

L'opzione \code{--prune} rimuove i pacchetti non più dichiarati nel file.

\manualnote[Ambiente validato]{
\filepath{environment.yml} fissa Python e lo stack Linux/CUDA usato nel
benchmark. Le build di PyTorch sono prelevate dall'indice ufficiale CUDA 12.6
di PyTorch, dichiarato nello stesso file.}

\subsection{Installazione del decoder Shurikode}

Il decoder, non incluso in \filepath{requirements.txt}, è disponibile nel
\href{https://github.com/Carraro-Riccardo/Shurikode-Context}{repository GitHub
Shurikode Context}. Dopo avere attivato \code{graph_rec}, clonarlo e installare
il pacchetto dalla sorgente:

\begin{lstlisting}[language={}]
git clone https://github.com/Carraro-Riccardo/Shurikode-Context.git
cd Shurikode-Context
python -m pip install .
\end{lstlisting}

I tre checkpoint dei modelli sono inclusi nel pacchetto e non richiedono un
download separato.

L'import atteso dal codice è:

\begin{lstlisting}[language={}]
from shurikodecontext import Dec
\end{lstlisting}

Se questo import non funziona, la pipeline non riuscirà a decodificare i simboli
anche se YOLO è installato correttamente.

\subsection{Installazione e runtime VGGT}

VGGT va installato dalla sorgente ufficiale
\href{https://github.com/facebookresearch/vggt}{Facebook Research}. Dopo avere
attivato \code{graph_rec}, installare la revisione usata nello sviluppo:

\begin{lstlisting}[language={}]
python -m pip install "git+https://github.com/facebookresearch/vggt.git@44b3afbd1869d8bde4894dd8ea1e293112dd5eba"
\end{lstlisting}

Il commit indicato nel comando rende riproducibile la versione di VGGT usata
dal progetto.

Al primo utilizzo, i pesi del modello configurato vengono scaricati
automaticamente e salvati nella cache locale. La prima esecuzione può quindi
richiedere una connessione Internet e risultare sensibilmente più lenta; le
esecuzioni successive riutilizzano i file già scaricati.

Nelle sessioni successive è sufficiente eseguire \cli{conda activate graph_rec}
dalla directory del repository.

\section{Requisiti hardware indicativi}

Il dimensionamento è stato ricavato da un'esecuzione completa senza cache sul
nodo \code{node9}, composto da sette viste. La tabella riporta sia i picchi
misurati sia la configurazione consigliata con un margine operativo.

\begin{table}[H]
\centering
\small
\caption{Sintesi del benchmark hardware}
\label{tab:benchmark-hardware}
\begin{tabularx}{\textwidth}{@{}>{\bfseries}p{0.34\textwidth}X@{}}
\toprule
Voce & Valore \\
\midrule
Scenario di prova & Cold run della pipeline su \code{node9}, con 7 viste \\
GPU di prova & NVIDIA RTX A5000 con 24~GiB di VRAM \\
Tempo di esecuzione & 81,9 secondi \\
Picco RAM del processo & 11,2~GiB \\
Picco VRAM misurato & 10,85~GiB \\
VRAM con margine del 20\% & 13,02~GiB \\
GPU consigliata & GPU NVIDIA compatibile con CUDA, con almeno 16~GiB di VRAM \\
RAM di sistema & 16~GiB minimi; 32~GiB consigliati \\
Stack validato & Python 3.10.21, PyTorch 2.6.0+cu126 e CUDA 12.6 \\
\bottomrule
\end{tabularx}
\end{table}

I valori non costituiscono requisiti assoluti: il consumo di memoria e i tempi
possono variare con il numero di viste, la risoluzione delle immagini, il
modello e la configurazione software utilizzata.
